\documentclass[11pt,a4paper,DIV9,german]{scrartcl}
%%%%%%%%%%%%%%%%%%%%%%%%%%%%%%%%%%%%%%%%%%%%%%%%%%%%%%%%%%%%%%%%%%%%%%%%%%
% Anpassungen f"ur das jeweilige "Ubungsblatt
\def\blnr{4}              % Blattnummer
\def\sdate{22.3.2019,}     % Abgabedatum
\def\exnr{1}              % Nummer der ersten Aufgabe des aktuellen Blattes
%%%%%%%%%%%%%%%%%%%%%%%%%%%%%%%%%%%%%%%%%%%%%%%%%%%%%%%%%%%%%%%%%%%%%%%%%%
\usepackage{babel,theorem}
\usepackage{exscale}
\usepackage{amsmath,amsfonts,amssymb,dsfont}
\usepackage{a4wide}
\usepackage{epic,eepic}
\usepackage{graphicx} % einbinden von Grafiken mit \includegraphics
\usepackage{url}

\pagestyle{empty}
\parindent0em
\parskip1ex
\addtolength{\topmargin}{0cm}
\addtolength{\oddsidemargin}{0cm}
\addtolength{\textheight}{2cm}
\addtolength{\textwidth}{0cm}

% neuer Z"ahler f"ur die Aufgaben
\newcounter{aufgnr}
\setcounter{aufgnr}{\exnr}
\addtocounter{aufgnr}{-1}

% in enumerate-Umgebungen erst Kleinbuchstaben, dann kleine r"omische Zahlen
\renewcommand{\labelenumi}{(\alph{enumi})\,}
\renewcommand{\labelenumii}{(\roman{enumii})\,}

% Stil f"ur Definitionen, S"atze usw.
\theoremstyle{break}   % Titelzeile abgesetzt
\newtheorem{Satz}{Satz}
\newtheorem{Lem}{Lemma}
\newtheorem{Def}{Definition}
\theorembodyfont{\upshape}  % normaler Font f"ur Beispiel, Aufgabe
\newtheorem{Bsp}{Beispiel}
\newtheorem{aufg}[aufgnr]{Aufgabe}
\newenvironment{Aufgabe}[1][z*+]%
{\begin{aufg}}
  {\end{aufg}\vspace{1.5ex}}
\newtheorem{praktaufg}[aufgnr]{Aufgabe^*}
\newenvironment{PraktAufgabe}[1][]%  der optionale Param. muss in [[ .. ]]
{\begin{praktaufg}#1}
  {\end{praktaufg}\vspace{1.5ex}}

\newcommand{\zweifaelles}[3]{ \left\{ \begin{array}{ccc} #1 & \hspace{.8cm}  \mbox{ falls } #2 \, , &     \\[1ex]
                               #3 & \mbox{ sonst } &     \end{array} \right.     }
% Mathematische K"urzel
\def\R{\mathbb{R}}
\def\S{\mathbb{S}}
\def\C{\mathbb{C}}
\def\Q{\mathbb{Q}}
\def\N{\mathbb{N}}
\def\Z{\mathbb{Z}}
\def\P{\mathbb{P}}
\def\K{\mathbb{K}}
\def\O{{\mathcal O} }
\def\G{{\mathcal G} }
\def\T{{\mathcal T} }
\def\U{{\mathcal U} }

\newcommand{\arrl}{\begin{array}{l}}
\newcommand{\are}{\end{array}}
\newcommand{\arrc}{\begin{array}{c}}
\newcommand{\mza}{\begin{eqnarray}}
\newcommand{\mze}{\end{eqnarray}}
\newcommand{\dn}{\frac{\partial}{\partial n}}
\newcommand{\dny}{\frac{\partial}{\partial n(y)}}
\newcommand{\dd}[2]{\ensuremath{\frac{\partial^{#1}}{\partial #2}}}
\newcommand{\ov}{\overline}
\newcommand{\ma}{\mid\!}
\newcommand{\me}{\!\mid}
\newcommand{\impl}{\Rightarrow}
\newcommand{\sca}[3]{\ensuremath{(#1,#2)_{#3}}}
\newcommand{\nor}[2]{\ensuremath{\parallel\! #1 \!\parallel_{#2}}}
\newcommand{\ra}{\rightarrow}
\newcommand{\phin}{\varphi}
\newcommand{\diver}{\mbox{div}}

\def\eps{{\varepsilon}}
\def\i{{\sl i\;}}
\def\e{{\rm e}}
\def\D{{\rm D\;}}
\def\norm{|\!|\!|}

\DeclareMathOperator{\spann}{span}
\DeclareMathOperator{\cond}{cond}
\DeclareMathOperator{\round}{round}
\DeclareMathOperator{\supp}{supp}

\begin{document}

\parbox{0ex}{ \scalebox{.27}{  \includegraphics[width=.5\textwidth,angle=-90]{cm-logo2.eps} }   }   \\

\parbox{25ex}{
  Prof.~Dr.~S.~Sauter\\
  Institut f"ur Mathematik\\
  Universit"at Z"urich
  }
%
\rule[0cm]{0.cm}{.01cm}
\hfill  \parbox{0.6\textwidth}{
  {\sf\LARGE\bfseries Numerik I}
  {\sf\Large\bfseries \;\;--\;\; Homework \blnr }\\[1.5ex]
  Deadline: \sdate\ 13:00
  }



\vspace{10ex}


% \scalebox{.27}{ \put(1250,0){ \includegraphics[width=.5\textwidth,angle=-90]{cm-logo2.eps} }  }


%%%%%%%%%%%%%%%%%%%%%%%%%%%%%%%%%%%%%%%%%%%%%%%%%%%%%%%%%%%%%%%%%%%%%%%%%%%
%
% Aufgaben
%
%%%%%%%%%%%%%%%%%%%%%%%%%%%%%%%%%%%%%%%%%%%%%%%%%%%%%%%%%%%%%%%%%%%%%%%%%%%

\textbf{Exercise 1} (8 P.) (Theoretical task)\\
Given a matrix \(\mathbf{A}_{\varepsilon}=\begin{bmatrix}\varepsilon&1\\1&1\end{bmatrix}\) for \(0\leq\varepsilon\leq \frac{1}{2}\).
\begin{enumerate}
  \item Compute the QR-decomposition \(\mathbf{A}_{\varepsilon}=\mathbf{Q}_{\varepsilon}\mathbf{R}_{\varepsilon}\) where \(\mathbf{Q}_{\varepsilon}\mathbf{Q}_{\varepsilon}^T=\mathbf{I}\) and \(\mathbf{R}_{\varepsilon}=\begin{bmatrix}r_{11}&r_{12}\\0&r_{22}\end{bmatrix}\).
  \item Compute the LU-decomposition \(\mathbf{A}_{\varepsilon}=\mathbf{L}_{\varepsilon}\mathbf{U}_{\varepsilon}\) where \(\mathbf{L}_{\varepsilon}=\begin{bmatrix}1&0\\l_{21}&1\end{bmatrix}\) and \(\mathbf{U}_{\varepsilon}=\begin{bmatrix}u_{11}&u_{12}\\0&u_{22}\end{bmatrix}\).
  \item The condition number of an \(n\times n\) matrix \(\mathbf{M}\) is defined as \(\kappa_*(\mathbf{M}):=\|\mathbf{M}\|_*\|\mathbf{M}^{-1}\|_*\) where \(\displaystyle \|\mathbf{M}\|_*:=\max_{1\leq{i,j}\leq n} |m_{ij}|\).  For a linear system \(\mathbf{M}\mathbf{u}=\mathbf{r}\), it indicates the influence of the  perturbation  in \(\tilde{\mathbf{r}}\approx \mathbf{r}\) on the perturbed solution of the system \(\mathbf{M}\tilde{\mathbf{u}}=\tilde{\mathbf{r}}\).
  \begin{enumerate}
    \item Calculate the condition numbers of \(\mathbf{A}_{\varepsilon}\), \(\mathbf{L}_{\varepsilon}\) and \(\mathbf{U}_{\varepsilon}\).  Investigate how the condition numbers change as \(\varepsilon\) approaches \(0\).
    \item Determine the condition numbers of \(\mathbf{Q}_\varepsilon\) and \(\mathbf{R}_\varepsilon\) for \(\varepsilon\rightarrow 0\).  (The norm is continuous, and therefore the norm and the limit are exchangeable.)
    \item Comment on your results in (i) and (ii).
  \end{enumerate}
\end{enumerate}

\vspace{0.6cm}

\textbf{Exercise 2} (6 P.) (Programming task)\\
Note: Do not use implemented Matlab functions to solve the exercises. 

a) Implement the QR-decomposition to solve a system of linear equations.  Write a function that takes the matrix \(\mathbf{A}\in\mathds{C}^{n\times n}\) as input  and returns the vectors \(\mathbf{u}, \; \beta\) as defined in \((3.5)\) of the lecture notes.


b) Write a program for solving $\mathbf{Ax}=\mathbf{b}$ with input $\mathbf{A}$ and the right-hand side vector \(\mathbf{b}\in\mathds{C}^{n\times 1}\) using your program from a) to get $\mathbf{u}$ and $\beta$. Use the algorithm in proof of Satz $3.6$ to efficiently solve the system. Adapt the algorithm on page $30$ to solve a system with upper-triangular matrix.  Use the following \(\mathbf{A}\) and \(\mathbf{b}\) to test your program: 
\begin{alignat*}{1}
  \mathbf{A}=&(a_{ij})_{i,j=1}^{100}\in\mathds{R}^{100\times 100}\quad\mathrm{where}~
  a_{ij}:=
  \begin{cases}
    100          ~&\mathrm{if}~i=j,\\
    e^{-\frac{|i-j|}{100}}~&\mathrm{if}~i\neq j.
  \end{cases}\\
  \mathbf{b}=&[1,0,1,0,\cdots,1,0]^T\in\mathds{R}^{100\times 1}.
\end{alignat*}

Check the result by calculating the residual \(\|\mathbf{r}\|_{\infty}=\|\mathbf{Ax}-\mathbf{b}\|_{\infty}\) where \(\displaystyle \|\mathbf{r}\|_{\infty}=\max_{1\leq i \leq n}|r_i|\).

\textbf{Exercise 3} (4 P.)(Theoretical task)\\
a) Prove \textbf{Lemma 3.5} of the lecture notes for the case $a_{11} = 0$.

b) Assume $dim(ker(\mathbf{A}))=m> 0.$ What is the consequence on the matrix $\mathbf{R}$ in (3.2) of the lecture notes?

\end{document}